\documentclass[a4paper,12pt]{article}
\usepackage[utf8]{inputenc}
\usepackage[english]{babel}
\usepackage[T1]{fontenc}
\usepackage{hyperref}
\usepackage{enumitem}
\usepackage{geometry} 
\usepackage{parskip} 
\usepackage[version=4]{mhchem}
\usepackage{amsmath}


\geometry{
	top=2.5cm,
	bottom=2.5cm,
	left=2.5cm,
	right=2.5cm
}


\hypersetup{
	colorlinks=true,
	linkcolor=blue,
	filecolor=magenta,      
	urlcolor=blue,
	pdftitle={Atmo-Ensemble Proposal},
}

\title{\textbf{Atmo-Ensemble}\\Ensemble ML Research Conclusion for Air Quality Multi-Classification}
\author{Josu Viteri \and Gotzon Viteri \and Iker Dominguez}
\date{March 2026}

\begin{document}
	
	\maketitle
	
	\begin{abstract}
		% Abstract
	\end{abstract}
	
	\section{Research Conclusions}
	
	% Purpose and expectatives
	
	\subsection{PCA and Dataset Feature Extraction}
		\textbf{Does studying correlated environmental features through principal component analysis (PCA) and other feature extraction techniques improve the robustness and efficiency of ensemble classifiers, or are non-linear feature interactions found by the models alone?}
		
			% Conclusion

	
	\subsection{Feature Importance and Interactions (XAI)}
	\label{subsec:xai}
	XAI context
	
	\subsection{Benchmark Performance}
	\textbf{Comparison of Bagging and Boosting}
	
	\textbf{Does the variance-reduction of parallel ensembles offer better generalization on noisy sensor data compared to the bias-reduction of sequential ensembles?}
	
	% Conclusion
	
	
	\subsection{Handling Ordinality}
	\textbf{Does encoding ordinality via Regressor-based Ensembles outperform standard Multiclass Classifiers in Air Quality Assessment?}

	% Conclusion
	
	
	\subsection{Imbalanced Learning}
	\subsubsection{Handling Rare Events}
	\textit{"Evaluating the efficacy of Hybrid Sampling-Ensemble methods (SMOTE-Boost vs. RUSBoost) in detecting rare Hazardous Air Quality events."}
	
	% Conclusion
	
	\subsubsection{Cost-Sensitive Learning vs. Resampling Methods}
	% Conclusion
	
	
	
	\subsubsection{Robustness to Missing Data (Simulating Broken Sensors)}
	% Conclusion
	
	\subsection{Semi-Supervised Learning}
	% Conclusion
	
	\subsection{Efficiency vs. Accuracy}
	\label{subsec:efficiency}
	\subsubsection{Trade-offs between Computational Cost and Classification Performance}
	\textit{Can lightweight ensembles replace deep stacking for edge-based Air Quality Monitoring?}
	
	% Conclusion
	
	
	\subsubsection{Conformal Prediction}
	% Conclusion
	
	\section{Limitations and Discovery}
	% Most common problems found

		\subsection{Limitations and Data Drift}
		% How the model might degrade over seasons or due to sensor aging.
	

	\section{Ethical and Social Implications}
	% How our problem helps society
		
		\subsection{Risk Assesment and Public Health
		% Impact of False Negatives in hazardous event detection.}

	
	\section{Summary and Next Steps}
	% Summary and next steps
	
\end{document}